\documentclass[10pt,a4paper]{report}

% Configuração do idioma e codificação
\usepackage[utf8]{inputenc}
\usepackage[T1]{fontenc}
\usepackage{newunicodechar}
\newunicodechar{–}{\textendash}
\usepackage[portuguese]{babel}

% Margens e formatação
\usepackage[a4paper,margin=1.7cm]{geometry}
\usepackage{graphicx}
\usepackage{amsmath}
\usepackage{booktabs}
\usepackage{array}
\usepackage{calc}
\usepackage{xurl}
\usepackage{microtype}
% Títulos de capítulos mais compactos.
\usepackage{titlesec}
\titleformat{\chapter}[display]
  {\normalfont\Large\bfseries}
  {\normalsize\chaptertitlename\ \thechapter}{8pt}{\Large}
\titleformat{name=\chapter,numberless}[display]
  {\normalfont\Large\bfseries}{}{0pt}{\Large}
\titlespacing*{\chapter}{0pt}{12pt}{20pt}
\usepackage{hyperref}
\graphicspath{{imagens/}}
\providecommand{\tightlist}{%
  \setlength{\itemsep}{0pt}\setlength{\parskip}{0pt}}
% Espaço adicional para a composição de linhas com títulos e URLs longos.
\setlength{\emergencystretch}{3em}
\usepackage{setspace}

\onehalfspacing
\setlength{\columnsep}{0.8cm}

% Informações do documento
\title{Sistema Inteligente de Apoio à Decisão
de Treinadores de Basquetebol}

\author{Relatório de AAUTIA e PLANTIA}

\date{\today}

\begin{document}

\maketitle

\begin{abstract}
Este trabalho investiga a integração de técnicas
de Machine Learning e Planeamento Automático
no apoio à decisão de treinadores de basquetebol.
\end{abstract}

\tableofcontents
\newpage

% Conteúdo dos capítulos em duas colunas.
\twocolumn
\input{capitulos/introducao}

\chapter{Estado da Arte}
\label{chap:estado-arte}
\input{capitulos/estado_arte_aautia}
\input{capitulos/estado_arte_plantia}

\input{capitulos/metodologia}
\input{capitulos/experiencias}
\input{capitulos/conclusoes}

\onecolumn
\phantomsection
\addcontentsline{toc}{chapter}{\bibname}
\bibliographystyle{plain}
\bibliography{referencias}

\end{document}
